\section{Introduction}

Object-goal navigation requires an agent to locate an instance of a target
category in an unknown environment, perceiving, deciding, and moving in a
closed loop \citep{batra2020objectnav}.  Recent zero-shot navigators delegate
the perceptual part of this loop to large vision--language models: a
panoramic observation is described in language, a heading is selected, and a
value map that scores candidate directions is updated at every step
\citep{long2024instructnav,yokoyama2024vlfm,zhou2023esc}. These methods perform
well on standard benchmarks, where the input is clean. In real deployments,
however, corridors are dimly lit, lenses fog up, and motion blurs the
image. Under these conditions, the same agents lose much of their success
rate: in our measurements, low light halves the success rate of a
representative zero-shot navigator on a held-out slice, and the paired
drop remains significant at $N{=}300$.

A natural remedy is to treat this as an image-quality issue and insert a
restoration front-end between the camera and the policy
\citep{guo2020zerodce,wu2023uretinex}, assuming that once the frame looks
clean  the navigator behaves as if the degradation were absent.
Whether this assumption holds has, to our knowledge, not been examined
directly for VLM-driven navigators. We find that it does not, and use a controlled
diagnosis to show \emph{why}. This is a diagnosis-and-benchmark paper: our
contribution is the measurement instrument and the resulting map of what
does not work, not a new agent.

\paragraph{A degradation-robust ObjectNav benchmark.}
We first build an evaluation substrate for this question: a degradation-robust
ObjectNav benchmark on a representative zero-shot backbone
(InstructNav~\citep{long2024instructnav}) in Habitat--Matterport~3D
\citep{ramakrishnan2021hm3d,yadav2023hm3dsem}. It applies structure-preserving
corruptions (low light, haze, motion blur, sensor noise) across severities,
fixes the same episodes across methods through a same-seed same-episode paired
protocol that supports McNemar tests, and equips the agent with a training-free
blind per-step reliability estimator. Every quantity below is measured on the
unmodified agent at the $N{=}300$ main-split scale, so the diagnosis does not
depend on any proposed fix. Sample size matters at this scale: between
$N{=}40$ and $N{=}150$ the measured per-condition restoration effect
\emph{changes sign}.

\paragraph{Degradation breaks commitments, not perception.}
Three observations emerge (Figure~\ref{fig:motivation}). First, degradation
reduces success: under low light the success rate falls by roughly half.
Second, the bottleneck is not appearance: a restoration step that raises image
fidelity by $8.4$\,dB PSNR leaves the success rate within seed noise, and a
per-frame oracle that picks the better of degraded and restored views does no
better than blanket restoration. Third, degradation does not change
\emph{whether} the agent acts. We call the moment the agent irreversibly
declares the target found---the detector's \texttt{found} flag flipping on a
single frame---a \emph{commitment}; degradation modestly lowers the commit rate
(from $93\%$ on clean to $89\%$ on low-light) and leaves
pure-exploration failures near  clean level, but corrupts the
\emph{quality} of these irreversible commitments. Decomposing the outcome as
$\text{SR}=\text{reachability}-\text{commit/stop loss}$, the agent walks within
one metre of the goal in $61\%$ of episodes yet succeeds in only $37\%$; it
commits at a frame reliability of $0.39$ versus $0.88$ on clean input, only
$42\%$ of its commitments lead to success, and once committed it rarely
recovers. The agent does not fail because it cannot see; it fails because
it makes an irreversible decision on a single unreliable frame and cannot
revoke it.

\paragraph{Negative results that localize the bottleneck.}
The diagnosis predicts that perception-layer fixes should not help, and we
confirm this at $N{=}300$: blanket restoration moves SR by only $+0.04$
(McNemar $p{=}0.18$), and even a non-deployable per-frame appearance
oracle caps at $+0.057$. A sweep across six causal families (appearance,
where-to-go, commit timing, forced commitment, temporal and multi-view
aggregation, and a depth-based cross-modal verifier) finds that every
\emph{deployable} family either lands at a non-significant ${\sim}0$ or
\emph{significantly hurts} success (Table~\ref{tab:sweep}), indicating a
commitment error that none of the probed test-time channels recovers.

\paragraph{Contributions.}
\begin{itemize}
\item \textbf{Benchmark.} A degradation-robust ObjectNav benchmark on a
  zero-shot VLM backbone: structure-preserving corruptions across
  severities, a same-seed paired protocol with McNemar exact tests and
  bootstrap CIs, and a training-free blind per-step reliability estimator.
\item \textbf{Diagnosis.} A method-independent mechanism analysis at
  $N{=}300$ that reframes degradation-induced failure as irreversible,
  low-reliability commitments, quantified by a reachability/commit-loss
  decomposition.
\item \textbf{Negative results.} A sweep of six causal intervention
  families in which every deployable family caps at zero or backfires
  ($p{<}0.05$), while the non-deployable per-frame appearance oracle caps
  at $+0.057$.
\item \textbf{Replication.} A cross-corruption replication (low-light,
  motion-blur, fog, Gaussian noise) reproducing the same non-significant
  fingerprint ($|\Delta\text{SR}|{\leq}0.05$, sign-flipping across
  cells), with every setting and seed released.
\end{itemize}
